% from ega1-4-fr.tex
\oldpage[I]{120}
Il suffit évidemment (\hyperref[I.2.1.3-fr]{2.1.3}) de considérer le cas où
$X$ est un schéma affine, et de montrer que dans ce cas $Y$ est fermé dans
$X$ et est un \emph{schéma affine}. En effet, si
$X=\operatorname{Spec}(A)$, on a $\mathscr{O}_X=\widetilde{A}$ et
$\mathscr{I}=\widetilde{\mathfrak{I}}$, où $\mathfrak{I}$ est un idéal de
$A$ (\hyperref[I.1.4.1-fr]{1.4.1})~; $Y$ est alors égal à la partie fermée
$V(\mathfrak{I})$ de $X$ et s'identifie au spectre premier de l'anneau
$B=A/\mathfrak{I}$ (\hyperref[I.1.1.11-fr]{1.1.11})~; de plus, si $\varphi$
est l'homomorphisme canonique $A\to B=A/\mathfrak{I}$, l'image directe
${}^a\varphi_*(\widetilde{B})$ s'identifie canoniquement au faisceau
$\widetilde{A}/\widetilde{\mathfrak{I}}=\mathscr{O}_X/\mathscr{I}$
(\hyperref[I.1.6.3-fr]{1.6.3} et \hyperref[I.1.3.9-fr]{1.3.9}), ce qui achève
la démonstration.

Nous dirons que $(Y,\mathscr{O}_Y)$ est le \emph{sous-préschéma} de
$(X,\mathscr{O}_X)$ \emph{défini par le faisceau d'idéaux $\mathscr{I}$}~;
c'est un cas particulier de la notion générale de \emph{sous-préschéma}~:

\begin{definition}[4.1.3]
\label{I.4.1.3-fr}
On dit qu'un espace annelé $(Y,\mathscr{O}_Y)$ est un
\emph{sous-préschéma} d'un préschéma $(X,\mathscr{O}_X)$ si~:
\begin{enumerate}
  \item[$1^\circ$] $Y$ est un sous-espace localement fermé de $X$~;
  \item[$2^\circ$] Si $U$ désigne le plus grand ouvert de $X$ contenant $Y$
  et tel que $Y$ soit fermé dans $U$ (autrement dit, le complémentaire dans
  $X$ de la frontière de $Y$ par rapport à $\overline{Y}$),
  $(Y,\mathscr{O}_Y)$ est un sous-préschéma de
  $(U,\mathscr{O}_X|U)$ défini par un faisceau quasi-cohérent d'idéaux de
  $\mathscr{O}_X|U$.
\end{enumerate}
On dit que le sous-préschéma $(Y,\mathscr{O}_Y)$ de $(X,\mathscr{O}_X)$ est
\emph{fermé} si $Y$ est fermé dans $X$ (auquel cas $U=X$).
\end{definition}

Il résulte aussitôt de cette définition et de
(\hyperref[I.4.1.2-fr]{4.1.2}) que les sous-préschémas fermés de $X$ sont en
\emph{correspondance biunivoque canonique} avec les faisceaux d'idéaux
quasi-cohérents $\mathscr{J}$ de $\mathscr{O}_X$, car le fait que deux tels
faisceaux $\mathscr{J}$, $\mathscr{J}'$ aient même support (fermé) $Y$ et
soient tels que les restrictions de $\mathscr{O}_X/\mathscr{J}$ et de
$\mathscr{O}_X/\mathscr{J}'$ à $Y$ soient identiques, entraîne aussitôt que
$\mathscr{J}'=\mathscr{J}$.

\begin{env}[4.1.4]
\label{I.4.1.4-fr}
Soient $(Y,\mathscr{O}_Y)$ un sous-préschéma de $X$, $U$ le plus grand ouvert
de $X$ contenant $Y$ et dans lequel $Y$ soit fermé, $V$ un ouvert de $X$
contenu dans $U$~; alors $V\cap Y$ est fermé dans $V$. En outre, si $Y$ est
défini par le faisceau quasi-cohérent $\mathscr{J}$ d'idéaux de
$\mathscr{O}_X|U$, $\mathscr{J}|V$ est un faisceau quasi-cohérent d'idéaux de
$\mathscr{O}_X|V$, et il est immédiat que le préschéma induit par $Y$ sur
$Y\cap V$ est le sous-préschéma fermé de $V$ défini par le faisceau d'idéaux
$\mathscr{J}|V$. Inversement~:
\end{env}

\begin{proposition}[4.1.5]
\label{I.4.1.5-fr}
Soit $(Y,\mathscr{O}_Y)$ un espace annelé tel que $Y$ soit un sous-espace de
$X$ et qu'il existe un recouvrement $(V_\alpha)$ de $Y$ par des ouverts de
$X$ tels que pour tout $\alpha$, $Y\cap V_\alpha$ soit fermé dans $V_\alpha$
et que l'espace annelé
$(Y\cap V_\alpha,\mathscr{O}_Y|(Y\cap V_\alpha))$ soit un sous-préschéma
fermé du préschéma induit sur $V_\alpha$ par $X$. Alors
$(Y,\mathscr{O}_Y)$ est un sous-préschéma de $X$.
\end{proposition}

L'hypothèse implique que $Y$ est localement fermé dans $X$ et que le plus
grand ouvert $U$ contenant $Y$ et dans lequel $Y$ est fermé contient tous les
$V_\alpha$~; on peut donc se ramener au cas où $U=X$ et $Y$ est fermé dans
$X$. On définit alors un faisceau quasi-cohérent d'idéaux $\mathscr{J}$ de
$\mathscr{O}_X$ en prenant pour $\mathscr{J}|V_\alpha$ le faisceau d'idéaux de
$\mathscr{O}_X|V_\alpha$ qui définit le sous-préschéma fermé
$(Y\cap V_\alpha,\mathscr{O}_Y|(Y\cap V_\alpha))$ et pour tout ouvert $W$ de
$X$ ne rencontrant pas $Y$, $\mathscr{J}|W=\mathscr{O}_X|W$. On vérifie
immédiatement en vertu de (\hyperref[I.4.1.3-fr]{4.1.3}) et
(\hyperref[I.4.1.4-fr]{4.1.4}) qu'il existe un faisceau d'idéaux $\mathscr{J}$
et un seul satisfaisant à ces conditions et qu'il définit le sous-préschéma
fermé $(Y,\mathscr{O}_Y)$.

En particulier, le préschéma \emph{induit} par $X$ sur un \emph{ouvert} de
$X$ est un \emph{sous-préschéma} de $X$.

\begin{proposition}[4.1.6]
\label{I.4.1.6-fr}
Un sous-préschéma (resp. un sous-préschéma fermé) d'un
sous-